\documentclass[11pt]{article}

\usepackage[T1]{fontenc}
\usepackage[utf8]{inputenc}
\usepackage{lmodern}
\usepackage{microtype}
\usepackage[margin=1in]{geometry}
\usepackage{amsmath,amssymb,amsthm,mathtools}
\usepackage{booktabs,array,longtable}
\usepackage{enumitem}
\usepackage{xcolor}
\usepackage{hyperref}
\usepackage{listings}
\usepackage{fancyhdr}

\hypersetup{
  colorlinks=true,
  linkcolor=blue!55!black,
  citecolor=blue!55!black,
  urlcolor=blue!55!black,
  pdftitle={Constructive divergence of the sum of reciprocal primes in iPRA},
  pdfauthor={OpenAI}
}

\pagestyle{fancy}
\setlength{\headheight}{14pt}
\fancyhf{}
\lhead{Prime reciprocal divergence in iPRA}
\rhead{Research note}
\cfoot{\thepage}

\newtheorem{theorem}{Theorem}[section]
\newtheorem{lemma}[theorem]{Lemma}
\newtheorem{corollary}[theorem]{Corollary}
\newtheorem{proposition}[theorem]{Proposition}
\theoremstyle{definition}
\newtheorem{definition}[theorem]{Definition}
\newtheorem{remark}[theorem]{Remark}

\newcommand{\N}{\mathbb{N}}
\newcommand{\Q}{\mathbb{Q}}
\newcommand{\ipra}{\mathrm{iPRA}}
\newcommand{\Prime}{\operatorname{Prime}}
\newcommand{\ProdP}{Q}
\newcommand{\NumP}{R}
\newcommand{\DenE}{D}
\newcommand{\Fact}{F}
\newcommand{\NumH}{T}
\newcommand{\Wit}{W}
\newcommand{\code}[1]{\texttt{#1}}

\lstdefinelanguage{Lean}{
  morekeywords={def,theorem,lemma,namespace,end,by,have,show,from,import,open,section,variable,where,if,then,else,let,in,match,with,case,induction,exact,refine,simp,rw,constructor},
  sensitive=true,
  morecomment=[l]{--},
  morecomment=[s]{/-}{-/},
  morestring=[b]"
}
\lstset{
  language=Lean,
  basicstyle=\ttfamily\small,
  columns=fullflexible,
  breaklines=true,
  frame=single,
  showstringspaces=false,
  backgroundcolor=\color{black!2}
}

\title{\textbf{Constructive Divergence of the Sum of Reciprocal Primes}\\
       \large A quantitative proof in intuitionistic primitive recursive arithmetic}
\author{Research note prepared for Vladimir Reshetnikov}
\date{15 November 2026}

\begin{document}
\maketitle

\begin{abstract}
We give a completely finite and constructive proof, formalizable in
intuitionistic primitive recursive arithmetic (\(\ipra\)), that the partial
sums of the reciprocals of the primes are unbounded.  The proof establishes
the explicit quantitative statement
\[
  \sum_{\substack{p\leq \Wit(k)\\p\ \mathrm{prime}}}\frac1p>k,
  \qquad
  \Wit(k)=2^{\,2\cdot 4^k}
  \quad(k\in\N).
\]
There are no limits, logarithms, real exponentials, infinite products,
impredicative definitions, choice principles, or uses of excluded middle for
non-decidable propositions.  The argument uses only finite products and sums,
a bounded factorization algorithm, binomial estimates, and quantifier-free
induction.  We also give a denominator-cleared internal formulation, a
detailed audit showing that every construction is primitive recursive, and a
Lean formalization plan tied to the current \code{ConciseIPRA} infrastructure
in the ProveIt repository.  The accompanying archive contains this source and
PDF, a Lean scaffold, an obligation-by-obligation formalization plan, and an
executable finite sanity checker.
\end{abstract}

\tableofcontents

\section{The constructive statement}

The usual analytic notation
\(
  \sum_p 1/p=+\infty
\)
is stronger-looking than what is actually needed.  Since every partial sum is
rational and the partial sums are increasing, the appropriate finitistic
statement is the following explicit unboundedness theorem.

\begin{theorem}[Quantitative prime-reciprocal divergence]\label{thm:main}
Let
\[
  \Wit(k):=2^{\,2\cdot 4^k}.
\]
For every natural number \(k\),
\[
  \sum_{\substack{p\leq \Wit(k)\\p\ \mathrm{prime}}}\frac1p>k.
\]
Moreover, the universal closure of a denominator-cleared version of this
inequality is provable in \(\ipra\).
\end{theorem}

The bound is intentionally elementary rather than sharp.  Its role is to make
the witness a single primitive recursive term and to make all proof
obligations quantifier-free after bounded searches have been compiled into
primitive recursive characteristic functions.

\subsection{Natural-number coding of the target}

For \(n\geq 2\), define
\begin{align*}
  \ProdP_n &:= \prod_{\substack{p\leq n\\p\ \mathrm{prime}}}p,\\
  \NumP_n &:= \sum_{\substack{p\leq n\\p\ \mathrm{prime}}}\frac{\ProdP_n}{p}.
\end{align*}
Every displayed quotient is exact.  Hence
\[
  \sum_{\substack{p\leq n\\p\ \mathrm{prime}}}\frac1p
  =\frac{\NumP_n}{\ProdP_n}.
\]
Consequently Theorem~\ref{thm:main} is exactly the natural-number assertion
\begin{equation}\label{eq:internal-target}
  \NumP_{\Wit(k)}>k\ProdP_{\Wit(k)}.
\end{equation}
No rational-number sort is required inside arithmetic.

The functions \(\ProdP\), \(\NumP\), and \(\Wit\) are primitive recursive.
For example, with \(\chi_{\Prime}(p)\in\{0,1\}\) the primitive recursive
primality characteristic,
\begin{align*}
 \ProdP(n)&=\prod_{p=2}^{n}
    \begin{cases}p,&\chi_{\Prime}(p)=1,\\1,&\chi_{\Prime}(p)=0,\end{cases}\\
 \NumP(n)&=\sum_{p=2}^{n}
    \begin{cases}\ProdP(n)\mathbin{\mathrm{div}}p,&\chi_{\Prime}(p)=1,\\0,&\chi_{\Prime}(p)=0.
    \end{cases}
\end{align*}
Bounded products, bounded sums, totalized division, exponentiation, and the
primality test are all primitive recursive.  Thus
\eqref{eq:internal-target} can be represented by a quantifier-free equation
between primitive recursive terms, for instance by applying a Boolean
less-than characteristic function and asserting that its value is \(1\).

\section{Finite auxiliary quantities}

We use two more finite rational quantities, always understood through their
natural numerators and denominators.

\begin{definition}[Harmonic numerator]\label{def:harmonic}
Let
\[
  \Fact_n:=n!,\qquad
  \NumH_n:=\sum_{m=1}^{n}\frac{\Fact_n}{m},\qquad
  H_n:=\frac{\NumH_n}{\Fact_n}=\sum_{m=1}^{n}\frac1m.
\]
Each \(m\leq n\) divides \(n!\), so the defining quotients are exact.
\end{definition}

\begin{definition}[Finite Euler product]\label{def:euler}
Let
\[
  \DenE_n:=\prod_{\substack{p\leq n\\p\ \mathrm{prime}}}(p-1),
  \qquad
  E_n:=\frac{\ProdP_n}{\DenE_n}
      =\prod_{\substack{p\leq n\\p\ \mathrm{prime}}}\frac{p}{p-1}.
\]
\end{definition}

All four functions \(\Fact,\NumH,\DenE,E\) are given by bounded primitive
recursion.  In the formal proof, statements about \(H_n\) and \(E_n\) are
replaced by cross-multiplied inequalities.  For example,
\[
 H_n\leq E_n
 \quad\Longleftrightarrow\quad
 \NumH_n\DenE_n\leq \Fact_n\ProdP_n.
\]

\section{A bounded constructive factorization lemma}

The Euler comparison below needs only finite factorization of the integers
\(1,\ldots,n\).  We record an explicitly primitive recursive construction so
that no appeal to an infinitary or classical form of unique factorization is
hidden in the proof.

\subsection{The factor trace}

For \(x>1\), let \(\ell(x)\) be the least divisor of \(x\) in the interval
\([2,x]\); set \(\ell(0)=\ell(1)=1\).  The function \(\ell\) is primitive
recursive: scan the bounded interval and retain the first divisor.  If
\(x>1\), then \(\ell(x)\) is prime.  Indeed, if
\(\ell(x)=ab\) with \(1<a<\ell(x)\), then \(a\mid x\), contradicting the
minimality of \(\ell(x)\).

Given \(m\), define a trace of length \(m+1\) by bounded iteration:
\[
 x_0=m,
 \qquad
 x_{i+1}=
 \begin{cases}
  1,&x_i\leq 1,\\
  x_i/\ell(x_i),&x_i>1.
 \end{cases}
\]
Whenever \(x_i>1\), one has \(x_{i+1}<x_i\).  Quantifier-free induction on
\(i\) proves that the trace reaches \(1\) by time \(m\).  Record each selected
prime \(\ell(x_i)\), and for every \(p\leq m\) let \(e_p(m)\) be the number of
occurrences of \(p\) in this record.  A second induction on the trace gives
\begin{equation}\label{eq:reconstruct}
  m=\prod_{\substack{p\leq m\\p\ \mathrm{prime}}}p^{e_p(m)}.
\end{equation}
All functions in this paragraph are primitive recursive.  A finite vector
\((e_2,\ldots,e_n)\) can, if desired, be encoded in base \(n+1\); digit
extraction is bounded division and remainder.

\begin{lemma}[Finite factor-vector injection]\label{lem:factor-injection}
Fix \(n\).  The map
\[
  m\longmapsto (e_p(m))_{\substack{p\leq n\\p\ \mathrm{prime}}}
  \qquad(1\leq m\leq n)
\]
is injective, every coordinate is at most \(n\), and
\[
  m=\prod_{\substack{p\leq n\\p\ \mathrm{prime}}}p^{e_p(m)}.
\]
\end{lemma}

\begin{proof}
For \(p>m\), the trace cannot select \(p\), so extending
\eqref{eq:reconstruct} from \(p\leq m\) to \(p\leq n\) only inserts factors
\(p^0=1\).  Each coordinate counts at most the \(m\) trace positions, hence is
at most \(n\).  If two integers give the same vector, reconstruction gives the
same product and therefore the integers are equal.  This proof is entirely
constructive; the only case distinctions concern decidable equality,
order, and divisibility of natural numbers.
\end{proof}

\begin{remark}[Why this is an \(\ipra\) proof]
The assertion that the trace has reached \(1\), the reconstruction invariant,
and the injectivity conclusion can all be written as equations between
primitive recursive terms.  The induction formulas are therefore
quantifier-free.  Course-of-values recursion is unnecessary: the trace is an
ordinary primitive recursion on a fuel parameter bounded by \(m\).
\end{remark}

\section{The harmonic sum is below a finite Euler product}

\begin{lemma}[Finite Euler domination]\label{lem:euler-domination}
For every \(n\geq 2\),
\[
  H_n\leq E_n.
\]
Equivalently,
\[
  \NumH_n\DenE_n\leq \Fact_n\ProdP_n.
\]
\end{lemma}

\begin{proof}
For each prime \(p\leq n\), consider the finite geometric sum
\[
  G_{p,n}:=1+\frac1p+\frac1{p^2}+\cdots+\frac1{p^n}.
\]
Expand the finite product
\[
  G_n:=\prod_{\substack{p\leq n\\p\ \mathrm{prime}}}G_{p,n}.
\]
Its terms are indexed by exponent vectors \((a_p)_p\) with
\(0\leq a_p\leq n\), and the term indexed by such a vector is
\(\prod_p p^{-a_p}\).  By Lemma~\ref{lem:factor-injection}, each
\(m\in\{1,\ldots,n\}\) determines a distinct exponent vector and contributes
the term
\[
  \prod_p p^{-e_p(m)}=\frac1m.
\]
All terms are nonnegative, so
\begin{equation}\label{eq:harmonic-truncated-euler}
  H_n\leq G_n.
\end{equation}

For each prime \(p\leq n\), finite geometric summation gives
\[
 G_{p,n}
 =\frac{p^{n+1}-1}{p^n(p-1)}
 \leq\frac{p^{n+1}}{p^n(p-1)}
 =\frac{p}{p-1}.
\]
Multiplying these finitely many inequalities and using
\eqref{eq:harmonic-truncated-euler} yields
\[
  H_n\leq
  \prod_{\substack{p\leq n\\p\ \mathrm{prime}}}\frac{p}{p-1}
  =E_n.
\]
Every displayed rational inequality is shorthand for a natural-number
inequality obtained by multiplying by a positive common denominator.
\end{proof}

\subsection{Internal proof shape}

For a literal \(\ipra\) derivation, the expansion can be organized by
induction on \(q=2,\ldots,n\), inserting the factor
\(1+p^{-1}+\cdots+p^{-n}\) when \(q=p\) is prime and inserting \(1\) when
\(q\) is composite.  Alternatively, encode exponent vectors in base
\(n+1\) and prove the finite distributive-law identity by induction on the
number of coordinates.  Both methods use only bounded sums/products and
quantifier-free induction.  No set-valued object is needed.

\section{A uniform bound for one Euler factor}

The next estimate replaces analytic inequalities such as
\(1+x\leq e^x\) by a finite binomial argument.

\begin{lemma}[Binomial bound]\label{lem:binomial-three}
For every integer \(a\geq 1\),
\[
  \left(1+\frac1a\right)^a\leq 3,
  \qquad\text{equivalently}\qquad
  (a+1)^a\leq 3a^a.
\]
\end{lemma}

\begin{proof}
By the binomial theorem,
\[
 \left(1+\frac1a\right)^a
 =\sum_{j=0}^{a}\binom aj\frac1{a^j}.
\]
For each \(j\),
\[
  \binom aj
  =\frac{a(a-1)\cdots(a-j+1)}{j!}
  \leq\frac{a^j}{j!},
\]
so the preceding sum is at most \(\sum_{j=0}^{a}1/j!\).  For \(j\geq1\),
quantifier-free induction gives \(j!\geq2^{j-1}\).  Hence
\[
 \sum_{j=0}^{a}\frac1{j!}
 \leq 1+\sum_{j=1}^{a}\frac1{2^{j-1}}
 < 1+2=3.
\]
Replacing the last strict inequality by a weak one proves the stated bound.
All denominators can be cleared; factorial, binomial coefficients, and the
finite geometric sum are primitive recursive.
\end{proof}

\begin{lemma}[Local Euler-factor bound]\label{lem:local-factor}
For every integer \(p\geq2\),
\begin{equation}\label{eq:local-factor-rational}
  \left(\frac{p}{p-1}\right)^p\leq4,
\end{equation}
or, equivalently,
\begin{equation}\label{eq:local-factor-integer}
  p^p\leq4(p-1)^p.
\end{equation}
\end{lemma}

\begin{proof}
For \(p=2\), equality holds.  For \(p=3\), the integer inequality is
\(27\leq32\).  Suppose \(p\geq4\), and put \(a=p-1\), so \(a\geq3\).
Lemma~\ref{lem:binomial-three} gives
\[
  (a+1)^a\leq3a^a.
\]
Multiplying by \(a+1\) and using \(3(a+1)\leq4a\), valid for \(a\geq3\),
we obtain
\[
  (a+1)^{a+1}\leq3a^a(a+1)\leq4a^{a+1}.
\]
Substituting \(p=a+1\) proves \eqref{eq:local-factor-integer}.
\end{proof}

\section{The Euler product is controlled by the prime reciprocal sum}

The crucial device is to raise the finite Euler product to the integer
\(\ProdP_n\).  This removes every fractional exponent that would appear in a
usual exponential estimate.

\begin{lemma}[Power-cleared product estimate]\label{lem:power-cleared}
For every \(n\geq2\),
\begin{equation}\label{eq:power-cleared-rational}
  E_n^{\ProdP_n}\leq4^{\NumP_n}.
\end{equation}
Equivalently, in natural-number form,
\begin{equation}\label{eq:power-cleared-integer}
  \ProdP_n^{\ProdP_n}
  \leq 4^{\NumP_n}\DenE_n^{\ProdP_n}.
\end{equation}
\end{lemma}

\begin{proof}
Every prime \(p\leq n\) divides \(\ProdP_n\).  Therefore
\begin{align*}
 E_n^{\ProdP_n}
 &=\prod_{\substack{p\leq n\\p\ \mathrm{prime}}}
      \left(\frac{p}{p-1}\right)^{\ProdP_n}\\
 &=\prod_{\substack{p\leq n\\p\ \mathrm{prime}}}
      \left[\left(\frac{p}{p-1}\right)^p\right]^{\ProdP_n/p}\\
 &\leq\prod_{\substack{p\leq n\\p\ \mathrm{prime}}}
      4^{\ProdP_n/p}\\
 &=4^{\sum_{p\leq n,\,p\ \mathrm{prime}}\ProdP_n/p}
 =4^{\NumP_n},
\end{align*}
where Lemma~\ref{lem:local-factor} supplies the inequality.  Clearing the
positive denominator \(\DenE_n^{\ProdP_n}\) gives
\eqref{eq:power-cleared-integer}.
\end{proof}

\begin{lemma}[A bounded prime sum bounds the Euler product]\label{lem:sum-bounds-product}
Let \(n\geq2\) and \(k\in\N\).  If
\[
  \NumP_n\leq k\ProdP_n,
\]
then
\begin{equation}\label{eq:Q-D-bound}
  \ProdP_n\leq4^k\DenE_n,
\end{equation}
and hence \(E_n\leq4^k\).
\end{lemma}

\begin{proof}
By Lemma~\ref{lem:power-cleared}, the hypothesis, and monotonicity of natural
exponentiation,
\[
 \ProdP_n^{\ProdP_n}
 \leq4^{\NumP_n}\DenE_n^{\ProdP_n}
 \leq4^{k\ProdP_n}\DenE_n^{\ProdP_n}
 =(4^k\DenE_n)^{\ProdP_n}.
\]
The exponent \(\ProdP_n\) is positive.  The map \(x\mapsto x^{\ProdP_n}\)
is strictly increasing on \(\N\), a fact proved by quantifier-free induction.
It follows that \(\ProdP_n\leq4^k\DenE_n\), which is
\eqref{eq:Q-D-bound}.  Dividing by the positive \(\DenE_n\) is merely
rational shorthand for the same cross-multiplied inequality.
\end{proof}

\begin{corollary}[Bounded prime sum bounds the harmonic sum]\label{cor:harmonic-upper}
If \(n\geq2\) and \(\NumP_n\leq k\ProdP_n\), then
\begin{equation}\label{eq:harmonic-upper-integer}
  \NumH_n\leq4^k\Fact_n,
\end{equation}
that is, \(H_n\leq4^k\).
\end{corollary}

\begin{proof}
Lemma~\ref{lem:euler-domination} gives
\(
 \NumH_n\DenE_n\leq\Fact_n\ProdP_n.
\)
Lemma~\ref{lem:sum-bounds-product} gives
\(
 \ProdP_n\leq4^k\DenE_n.
\)
Thus
\[
 \NumH_n\DenE_n
 \leq\Fact_n\ProdP_n
 \leq4^k\Fact_n\DenE_n.
\]
Since \(\DenE_n>0\), cancellation yields
\eqref{eq:harmonic-upper-integer}.
\end{proof}

\section{A dyadic lower bound for the harmonic sum}

\begin{lemma}[Dyadic harmonic lower bound]\label{lem:dyadic}
For every \(t\in\N\),
\begin{equation}\label{eq:dyadic-rational}
  H_{2^{2t}}\geq1+t.
\end{equation}
Equivalently,
\begin{equation}\label{eq:dyadic-integer}
  \NumH_{2^{2t}}\geq(1+t)\Fact_{2^{2t}}.
\end{equation}
\end{lemma}

\begin{proof}
For an arbitrary \(r\), partition the terms of \(H_{2^r}\) after the initial
term \(1\) into the \(r\) blocks
\[
  2^j<m\leq2^{j+1}\qquad(0\leq j<r).
\]
The \(j\)-th block contains \(2^j\) terms.  Every denominator in it is at most
\(2^{j+1}\), so every reciprocal is at least \(2^{-(j+1)}\).  Therefore each
block has sum at least
\[
  2^j\cdot2^{-(j+1)}=\frac12.
\]
Consequently
\[
  H_{2^r}\geq1+\frac r2.
\]
Taking \(r=2t\) gives \eqref{eq:dyadic-rational}.  Multiplication by the
positive factorial denominator gives \eqref{eq:dyadic-integer}.  Formally, the
block estimate is a bounded-sum inequality, and the accumulation of the
blocks is proved by quantifier-free induction on \(j\).
\end{proof}

\section{Proof of the main theorem}

\begin{proof}[Proof of Theorem~\ref{thm:main}]
Fix \(k\in\N\), put
\[
  t:=4^k,
  \qquad
  n:=2^{2t}=\Wit(k).
\]
Assume, for contradiction, that the prime reciprocal partial sum is at most
\(k\).  In denominator-cleared form this is
\begin{equation}\label{eq:contra-assumption}
  \NumP_n\leq k\ProdP_n.
\end{equation}
By Corollary~\ref{cor:harmonic-upper},
\begin{equation}\label{eq:upper-final}
  \NumH_n\leq4^k\Fact_n=t\Fact_n.
\end{equation}
On the other hand, Lemma~\ref{lem:dyadic} gives
\begin{equation}\label{eq:lower-final}
  \NumH_n\geq(1+t)\Fact_n>t\Fact_n,
\end{equation}
contradicting \eqref{eq:upper-final}.  Hence
\(
 \neg(\NumP_n\leq k\ProdP_n).
\)
Order on natural numbers is decidable, so intuitionistically this implies
\(
 \NumP_n>k\ProdP_n.
\)
Equivalently,
\[
  \sum_{\substack{p\leq n\\p\ \mathrm{prime}}}\frac1p>k.
\]
The witness \(n=\Wit(k)\) is primitive recursive and was fixed before the
contradiction argument; no choice principle is used.
\end{proof}

\begin{corollary}[Constructive divergence]\label{cor:divergence}
The monotone sequence of rational partial sums
\(
 S_n=\sum_{p\leq n}1/p
\)
diverges to \(+\infty\) in the constructive sense
\[
  \forall k\in\N\;\exists n\in\N\;(S_n>k),
\]
and in fact the witness can be chosen as \(n=\Wit(k)\).
\end{corollary}

For a rational bound \(a/b\), choose any natural \(k>a/b\), for example
\(k=a+1\).  Then the same explicit partial sum exceeds \(a/b\).  Thus the
natural-bound formulation loses no ordinary content.

\section{Formalizability in intuitionistic PRA}

This section audits the proof against the restrictions of \(\ipra\).  The
point is not merely that every computation is finite; the proof must be
arranged so that every induction instance has a quantifier-free matrix.

\subsection{The presentation of iPRA}

The ProveIt \code{ConciseIPRA} development contains the usual first-order
presentation \code{PRAAx}: one function symbol for every primitive recursive
function, defining equations for zero, successor, projection, composition,
and primitive recursion, successor separation and injectivity, and induction
for quantifier-free formulas.  Its theorem \code{PRA.qfInd} exposes the
quantifier-free induction rule.  The file \code{PRAToScalar.lean} defines a
translation \code{FromPRA.code} of every primitive recursive symbol to the
scalar iterator language and proves that every \code{PRAAx} axiom is a theorem
of the concise scalar theory.  The general preservation theorem is
\code{Interp.sound}.

Accordingly, a clean formalization has two layers:
\begin{enumerate}[label=\arabic*.]
\item derive the denominator-cleared theorem in \code{PRAAx};
\item apply \code{FromPRA.interp.sound} to obtain its translation in the
      concise scalar core.
\end{enumerate}
This avoids reproving elementary number theory directly in the smallest
seven-, six-, or five-rule presentation.

\subsection{Primitive recursive inventory}

The following functions are needed, all by ordinary primitive recursion and
composition:
\begin{center}
\begin{tabular}{p{0.29\textwidth}p{0.62\textwidth}}
\toprule
Function or predicate & Primitive recursive construction \\
\midrule
Equality, order, sign, conditional & Standard recursion on natural numbers. \\
Divisibility and remainder & Bounded subtraction/scan, or the usual primitive recursive division algorithm. \\
Primality characteristic & Check \(2\leq p\) and scan the bounded divisors \(2,\ldots,p-1\). \\
Bounded sum and product & Primitive recursion on the upper bound. \\
Power, factorial, binomial & Primitive recursion; binomial may use Pascal recursion or factorial identities with exact division. \\
Least divisor \(\ell(x)\) & Bounded scan retaining the first divisor. \\
Factor trace and exponent counts & Primitive recursion on a fuel parameter; each update is bounded arithmetic. \\
Finite-vector coding & Base-\((n+1)\) coding, with bounded quotient and remainder for digits. \\
\(\ProdP,\NumP,\DenE,\Fact,\NumH,\Wit\) & Composition of the preceding bounded folds and arithmetic functions. \\
\bottomrule
\end{tabular}
\end{center}

Bounded logical quantifiers occurring in an informal definition are not left
as quantifiers in an induction formula.  If \(b(i)\in\{0,1\}\) is a
characteristic function, then
\[
  \forall i\leq n\;P(i)
  \quad\text{is represented by}\quad
  \prod_{i=0}^{n}b(i)=1,
\]
and bounded existence can be represented by the sign of
\(\sum_{i=0}^{n}b(i)\).  Therefore every loop invariant used above has a
quantifier-free representation.

\subsection{Induction audit}

\begin{longtable}{p{0.24\textwidth}p{0.29\textwidth}p{0.38\textwidth}}
\toprule
Proof component & Induction parameter & Quantifier-free invariant \\
\midrule
\endfirsthead
\toprule
Proof component & Induction parameter & Quantifier-free invariant \\
\midrule
\endhead
Least-divisor trace terminates & Fuel \(i\leq m\) & Current value decreases when above \(1\); bounded comparisons are Boolean PR terms. \\
Factor reconstruction & Trace length & Accumulated selected-prime product times current quotient equals \(m\). \\
Finite-product expansion & Number of prime coordinates & Equality of denominator-cleared bounded sums indexed by base-coded vectors. \\
Geometric-factor estimate & Exponent bound & Standard finite geometric identity after cross multiplication. \\
Factorial lower bound & \(j\) & \(2^{j-1}\leq j!\) for \(j\geq1\), written with a conditional at \(j=0\). \\
Binomial theorem and estimate & \(a\) or binomial row index & Polynomial identity and termwise integer inequalities. \\
Power monotonicity & Positive exponent & \(x<y\Rightarrow x^r<y^r\), with order represented by a Boolean term. \\
Dyadic block estimate & Block index & Denominator-cleared lower bound for the accumulated bounded sum. \\
Main theorem & No induction needed at the final step & A decidable inequality contradiction with explicit \(n=\Wit(k)\). \\
\bottomrule
\end{longtable}

\subsection{Intuitionistic audit}

Only decidable numerical propositions are split into cases: equality, order,
divisibility, primality, and membership in explicitly bounded ranges.  The
last step uses
\[
 \neg(a\leq b)\Longrightarrow b<a,
\]
which is intuitionistically valid because natural order is decidable.  The
factorization witness is computed by a fixed bounded program.  No arbitrary
prime factor, subsequence, limit, or infinite product is selected.  Thus the
argument does not conceal an appeal to excluded middle or choice.

\section{Exact internal theorem for Lean}

Let \code{ltB : PR 2} be a primitive recursive Boolean less-than symbol, and
let \code{qPR}, \code{rPR}, and \code{wPR} be the symbols representing
\(\ProdP\), \(\NumP\), and \(\Wit\).  With \(x\) the sole free variable, the
quantifier-free target body is schematically
\[
 \varphi(x)\;:\equiv\;
 \operatorname{ltB}\bigl(x\cdot \ProdP(\Wit(x)),\,
                          \NumP(\Wit(x))\bigr)=1.
\]
The desired standard-iPRA theorem has the Lean type
\begin{lstlisting}
[] \vdash[PRAAx] Fm.all phi
\end{lstlisting}
where \code{phi : Fm PTm 1} and \code{phi.QF}.  Once this is established, the
concise-core theorem is obtained with the existing interpretation:
\begin{lstlisting}
FromPRA.interp.sound primeReciprocalDivergencePRA
\end{lstlisting}
Its conclusion maps every primitive recursive term through the morphism
\code{FromPRA.morph.f}.  The compiler \code{FromPRA.code} sends composition to
substitution and primitive recursion to the derived scalar recursion binder.

\section{Lean formalization roadmap}

The archive includes \code{lean/PrimeReciprocalIPRA.lean}, a theorem-signature
scaffold.  A complete formalization can be built in the following modules.

\subsection{Phase I: semantic finite mathematics in Mathlib}

\begin{enumerate}[label=I.\arabic*.]
\item \textbf{Definitions.}  Define the six finite functions listed in the
      accompanying scaffold: the prime set, its product and reciprocal-sum
      numerator, the Euler denominator, the harmonic numerator, and the
      witness.
\item \textbf{Factorization injection.}  For the first semantic proof, use
      Mathlib's finite factorization API or an explicit least-factor trace.
      State an injective map from \(m\leq n\) to bounded exponent vectors.
\item \textbf{Euler domination.}  Prove \(H_n\leq E_n\), then prove the
      cross-multiplied integer theorem used by the internal layer.
\item \textbf{Local factor.}  Prove
      \(p^p\leq4(p-1)^p\) for \(2\leq p\), with the binomial lemma separated
      out.
\item \textbf{Product control.}  Prove
      \[
      \ProdP_n^{\ProdP_n}\leq
      4^{\NumP_n}\DenE_n^{\ProdP_n}.
      \]
\item \textbf{Dyadic harmonic bound.}  Prove the integer statement
      \[
        (1+t)(2^{2t})!\leq \NumH_{2^{2t}}.
      \]
\item \textbf{Semantic main theorem.}  Combine the integer inequalities to
      obtain \(k\ProdP_{\Wit(k)}<\NumP_{\Wit(k)}\), and then derive the rational
      corollary.
\end{enumerate}

\subsection{Phase II: primitive recursive codes}

Build a small reusable library of \code{PR} terms for Boolean arithmetic,
bounded folds, quotient/remainder, and base-coded finite vectors.  Each
semantic function from Phase I receives:
\begin{enumerate}[label=II.\arabic*.]
\item a symbol \code{fPR : PR arity};
\item an evaluation theorem identifying its standard-natural semantics;
\item its defining equations as derivations from \code{PRAAx}.
\end{enumerate}
Because \code{PRAAx} already contains the defining equations for every
\code{PR} symbol, the main work is a library of derived equational and order
lemmas, not re-establishing computability.

\subsection{Phase III: syntactic iPRA derivation}

Mirror the mathematical dependency graph:
\begin{lstlisting}
factorTrace_correct
  -> factorVector_injective
  -> harmonic_le_euler_cross

factorial_ge_pow_two
  -> binomial_three_bound
  -> local_factor_bound
  -> euler_power_bound

harmonic_le_euler_cross + euler_power_bound
  -> harmonic_upper_of_prime_sum_upper

dyadic_harmonic_lower + harmonic_upper_of_prime_sum_upper
  -> primeReciprocalDivergencePRA
\end{lstlisting}
Every induction call should go through \code{PRA.qfInd}; the proof should not
use a semantic soundness argument to jump from truth in \(\N\) to derivability.
The semantic layer is a specification and debugging aid, while Phase III is
the actual formal proof in \(\ipra\).

\subsection{Phase IV: compilation to ConciseIPRA}

Apply \code{FromPRA.interp.sound}.  The current repository already proves that
projection and composition translate to reflexivity, primitive recursion to
the scalar recursion equations, and quantifier-free induction to the derived
scalar induction theorem.  Therefore no new metatheory is required for the
final transfer.

\section{Potential simplifications and improvements}

\begin{enumerate}
\item The witness \(2^{2\cdot4^k}\) is chosen for proof transparency.  A
      better local constant than \(4\), or a sharper harmonic block estimate,
      immediately improves it.
\item One may replace the full factor-vector expansion by an induction over
      \(q\leq n\) that groups integers according to their \(q\)-adic
      valuation.  This may shorten the syntactic Lean proof because it avoids
      an explicit finite-vector injection theorem.
\item A direct proof in the scalar iterator calculus could share iteration
      terms between primality testing, bounded products, and the factor trace.
      That is likely smaller as a certificate but considerably harder to
      maintain than the \code{PRAAx}-then-interpret route.
\item The denominator \(n!\) is deliberately simple.  Replacing it by
      \(\operatorname{lcm}(1,\ldots,n)\) reduces numeral growth but introduces
      gcd/lcm lemmas.  It is unlikely to reduce proof complexity in the first
      formalization.
\end{enumerate}

\section{Further research questions}

\begin{enumerate}
\item What is the smallest elementary witness function \(W(k)\) obtainable
      from a proof using only bounded products, qf induction, and no analytic
      functions?
\item Can the entire proof be organized around a single valuation recursion,
      eliminating the separately coded factor-vector expansion?
\item How small is the resulting proof certificate after translation to the
      current five- or six-rule ConciseIPRA kernels?
\item Is there a materially shorter proof in \(\ipra\) based on the Erdos
      counting argument, once convergence is replaced by an explicit bounded
      partial-sum hypothesis?
\item Can the method be generalized in \(\ipra\) to divergent reciprocal sums
      over primes in selected arithmetic progressions without importing
      analytic number theory?
\end{enumerate}

\section*{Repository and source notes}
\addcontentsline{toc}{section}{Repository and source notes}

The target theorem was taken from the expository page ``Divergence of the sum
of the reciprocals of the primes.''  The formalization plan was matched to the
ProveIt repository on its \code{main} branch at commit
\code{0e834f7851375f4b22ae8892000495cb6fa9eee8}, dated 15 November 2026.
In particular, the plan uses:
\begin{itemize}
\item \code{Logic/PrimitiveRecursiveArithmetic/ConciseIPRA/Lean/ConciseIPRA/PR.lean};
\item \code{.../PRAToScalar.lean};
\item \code{.../Interpretation.lean};
\item the root \code{lakefile.toml}, which registers the \code{ConciseIPRA}
      Lean library and pins Mathlib \code{v4.32.0}.
\end{itemize}

\begin{thebibliography}{9}
\bibitem{wiki}
Wikipedia contributors,
\emph{Divergence of the sum of the reciprocals of the primes},
\url{https://en.wikipedia.org/wiki/Divergence_of_the_sum_of_the_reciprocals_of_the_primes},
accessed 15 November 2026.

\bibitem{concise}
V. Reshetnikov,
\emph{Concise presentations of intuitionistic primitive recursive arithmetic},
ProveIt repository,
\url{https://github.com/VladimirReshetnikov/ProveIt/tree/main/Logic/PrimitiveRecursiveArithmetic/ConciseIPRA},
repository snapshot noted above.

\bibitem{prlean}
V. Reshetnikov,
\code{ConciseIPRA/PR.lean}, standard first-order presentation of \(\ipra\),
ProveIt repository.

\bibitem{pratoscalar}
V. Reshetnikov,
\code{ConciseIPRA/PRAToScalar.lean}, interpretation of standard \(\ipra\) in
the scalar theory, ProveIt repository.

\bibitem{interpretation}
V. Reshetnikov,
\code{ConciseIPRA/Interpretation.lean}, derivability preservation for term
interpretations, ProveIt repository.
\end{thebibliography}

\appendix

\section{Denominator-cleared dependency list}

For a kernel-oriented implementation, the complete numerical core can be
stated without rational notation as follows.  For \(n\geq2\):
\begin{align*}
\text{(E1)}\quad& \NumH_n\DenE_n\leq\Fact_n\ProdP_n,\\
\text{(E2)}\quad& p^p\leq4(p-1)^p\quad(2\leq p),\\
\text{(E3)}\quad& \ProdP_n^{\ProdP_n}
                   \leq4^{\NumP_n}\DenE_n^{\ProdP_n},\\
\text{(E4)}\quad& \NumP_n\leq k\ProdP_n
                   \Longrightarrow \ProdP_n\leq4^k\DenE_n,\\
\text{(E5)}\quad& \NumP_n\leq k\ProdP_n
                   \Longrightarrow \NumH_n\leq4^k\Fact_n,\\
\text{(E6)}\quad& \NumH_{2^{2t}}\geq(1+t)\Fact_{2^{2t}}.
\end{align*}
Substitute \(t=4^k\), \(n=2^{2t}\).  If
\(\NumP_n\leq k\ProdP_n\), (E5) gives
\(\NumH_n\leq t\Fact_n\), while (E6) gives
\(\NumH_n\geq(1+t)\Fact_n>t\Fact_n\).  Decidable order yields
\(\NumP_n>k\ProdP_n\).

\section{Suggested Lean theorem signatures}

The accompanying scaffold contains these signatures in context.  The central
integer theorems should look essentially as follows:
\begin{lstlisting}
theorem localFactorBound (p : Nat) (hp : 2 <= p) :
    p ^ p <= 4 * (p - 1) ^ p

theorem harmonicLeEulerCross (n : Nat) :
    harmonicNumerator n * eulerDenominator n <=
      Nat.factorial n * primeProduct n

theorem eulerPowerBound (n : Nat) :
    primeProduct n ^ primeProduct n <=
      4 ^ primeNumerator n *
        eulerDenominator n ^ primeProduct n

theorem dyadicHarmonicLower (t : Nat) :
    (1 + t) * Nat.factorial (2 ^ (2 * t)) <=
      harmonicNumerator (2 ^ (2 * t))

theorem quantitativeDivergenceNat (k : Nat) :
    k * primeProduct (witness k) <
      primeNumerator (witness k)
\end{lstlisting}

\section{What the archive does and does not claim}

The paper supplies a fully developed human-readable proof and a detailed
formalizability argument.  The Lean file is an outline with explicit theorem
signatures and \code{sorry} placeholders; it is not represented as a completed
kernel-checked formalization.  The Python checker is only a finite sanity
check of the arithmetic identities and inequalities on sample ranges.  The
proof of Theorem~\ref{thm:main} is mathematical and does not depend on those
computations.

\end{document}
